% Context from chapters/sec-autodiff.tex; for mathematical reference, not compilation.
erivative of the output node by
\eq{
	\pd{x_t}{x_t} = \Id_{n_t \times n_t}, 
}
and apply the following recursion from the last node to the first:
\eq{
	\foralls k=t-1,t-2,\ldots,1, \quad
	\pd{x_t}{x_k} 
	= \sum_{m \in \text{son}(k)}  \pd{x_t}{x_m} \times \bk{\pd{x_m}{x_k}}
	= \sum_{m \in \text{son}(k)} \pd{x_t}{x_m} \times \pd{f_m(x_1,\ldots,x_{m-1})}{x_k}.
}
 
The set $\text{son}(k)$ contains the children $m>k$ that depend directly on $x_k$, as shown in Figure~\ref{fig-forward-backward}, right. Contributions from every child must be added.

   
\paragraph{Back-propagation.}

For a scalar output $x_t \in \RR$, we have $\pd{x_t}{x_k} = [\nabla_{x_k} f(x)]^\top \in \RR^{1 \times n_k}$. The reverse recursion can then be written for gradient vectors:
\eq{
	\foralls k=t-1,t-2,\ldots,1, \quad
	\nabla_{x_k} f(x) 
	= \sum_{m \in \text{son}(k)} \pa{\pd{f_m(x_1,\ldots,x_{m-1})}{x_k}}^\top \pa{
	 	\nabla_{x_m} f(x)
	}.
}
Here $\pa{\pd{f_m(x_1,\ldots,x_{m-1})}{x_k}}^\top \in \RR^{n_k \times n_m}$ is the adjoint Jacobian of $f_m$. This adjoint recursion is called backpropagation and is the standard form used in machine learning.

For a scalar output ($n_t=1$), reverse mode is usually the preferred way to compute the full gradient. It requires access to the intermediate variables $(x_k)_{k=s+1}^t$, which can be costly to store for a large graph. Checkpointing trades additional computation for reduced storage.

   
\paragraph{Simple example.}

For the two-input example $f(x,y)$ in~\eqref{eq-simple-func-autodiff}, reverse mode gives
\begin{align*}
		\pd{f}{f} &= 1 &\\
		\pd{f}{c} &= \pd{f}{f} \pdb{f}{c} = \pd{f}{f} 1 &
			\{ c \mapsto f = b+c\}\\
		\pd{f}{b} &= \pd{f}{c} \pdb{c}{b} + \pd{f}{f} \pdb{f}{b} = \pd{f}{c} \frac{1}{2\sqrt{b}} + \pd{f}{f} 1 & 
			\{ b \mapsto c=\sqrt{b}, b \mapsto f=b+c\} \\
		\pd{f}{a} &= \pd{f}{b} \pdb{b}{a} = \pd{f}{b} y & 
			\{ a \mapsto b=y a\} \\
		\pd{f}{y} &= \pd{f}{b} \pdb{b}{y} = \pd{f}{b} a & 
			\{ y \mapsto b=ya\} \\
		\pd{f}{x} &= \pd{f}{a} \pdb{a}{x} = \pd{f}{a} \frac{1}{x} & 
			\{ x \mapsto a=\log(x)\}
\end{align*}
A single reverse traversal computes both derivatives with respect to $x$ and $y$, whereas coordinatewise forward mode requires two passes.

     
   
                                 

                                                                                                                   
              
                    
                                    
                     
                          
                                    
                              
                                                
                                                    
                             
                                                           
                                              
                                                       
                                                           
                                                        
                                                         
                                                        
                                                                
                                                            
                                                        
                                              
                                               
            
   


                                                                        
\subsection{Feed-forward Compositions}\label{sec-reversemode-simple}

For a feedforward chain, the computed function is the composition
\eql{\label{eq-simple-lin-dag}
	f = f_{t} \circ f_{t-1} \circ \ldots \circ f_2 \circ f_1
}
for functions $f_k : \RR^{n_{k-1}} \rightarrow \RR^{n_{k}}$.

The forward evaluation starts from $x_0=x \in \RR^{n_0}$ and computes
\eq{
	\foralls k=1, \ldots, t, \quad x_{k} = f_k(x_{k-1})
}
giving the output $f(x) = x_t$.

Let $A_k \eqdef \partial f_k(x_{k-1}) \in \RR^{n_{k} \times n_{k-1}}$ be the Jacobian of the local map. The chain rule gives
\eq{
	\partial f(x) = A_{t} \times A_{t-1} \times \ldots A_2 \times A_1.
}
Forward mode multiplies the Jacobians from right to left, while reverse mode multiplies them from left to right
\begin{align*}
	\partial f(x) &= A_{t} \times \pa{  A_{t-1} \times \pa{ \ldots \times \pa{ A_3 \times \pa{ A_2 \times A_1 } } } }, \\
	\partial f(x) &= \pa{ \pa{ \pa{ \pa{ A_{t} \times A_{t-1} }  \times A_{t-2} } \times \ldots } \times A_2 } \times A_1.
\end{align*}


\begin{figure}
\centering
\includegraphics[width=.42\linewidth]{auto-diff/matrix-forward} 
\qquad
\includegraphics[width=.45\linewidth]{auto-diff/matrix-backward} 
\caption{\label{fig-matrix-mult}
Matrix multiplication costs for forward mode (left) and reverse mode (right) along a chain.
}
\end{figure}

We note that the computation of the product $A \times B$ of $A \in \RR^{n \times p}$ with $B \in \RR^{p \times q}$ requires $O(npq)$ arithmetic operations with the standard dense algorithm.
 
As shown on Figure~\ref{fig-matrix-mult}, the arithmetic costs of the two parenthesizations are proportional to
\eq{
	n_0 \sum_{k=1}^{t-1} n_k n_{k+1}
	\qandq
	n_t \sum_{k=0}^{t-2} n_k n_{k+1}
}
For roughly comparable intermediate widths, these costs favor reverse mode when $n_t\ll n_0$. A scalar learning objective typically has $n_t=1$, making reverse mode a natural choice; the exact cost comparison depends on all the intermediate dimensions.


                                                                        
\subsection{Feed-forward Architecture}

A parameterized feedforward architecture, such as a neural network, has the form
\eql{\label{eq-feednets}
	\foralls k=1, \ldots, t, \quad x_{k} = f_k(x_{k-1},\th_{k-1})
}
where $\th_{k-1}$ is the parameter vector for a layer and the input $x_0 \in \RR^{n_0}$ is fixed. The objective is
\eql{\label{eq-loss-feedf}
	f(\th) \eqdef L(x_t)
}
Here $L : \RR^{n_t} \rightarrow \RR$ is a loss, such as squared prediction error or logistic loss, and $\th=(\th_k)_{k=0}^{t-1}$ collects the parameters. Figure~\ref{fig-mlp} shows the computational graph.


\begin{figure}
\centering
\includegraphics[width=.7\linewidth]{auto-diff/archi-feedforward} 
\caption{\label{fig-mlp}
Computational graph for a feedforward architecture.
}
\end{figure}


Reverse mode computes the gradient of $f$ by propagating derivatives through the variables $(x_k,\th_k)$. Initialize
\eq{
	\nabla_{x_t} f = \nabla L(x_t) 	
}
and recurse from $k=t$ down to $1$:
\eql{\label{eq-backprop-discr}
	z_{k-1} = [\partial_x f_k(x_{k-1},\th_{k-1})]^\top z_k
	\qandq
	\nabla_{\th_{k-1}} f = [\partial_\th f_k(x_{k-1},\th_{k-1})]^\top (\nabla_{x_{k}} f)
}
where $z_k \eqdef \nabla_{x_{k}} f(\th)$ is the gradient with respect to $x_{k}$.

   
\paragraph{Multilayer perceptrons.}\label{sec-a